
\documentclass[12pt]{article}

%==================== VIETNAMESE ====================
\usepackage[utf8]{inputenc}
\usepackage[T5]{fontenc}
\usepackage[vietnamese]{babel}

%==================== MATH ====================
\usepackage{amsmath, amssymb, amsthm}

%==================== ALGORITHM ====================
\usepackage{algorithm}
\usepackage{algpseudocode}

%==================== GRAPHICS ====================
\usepackage{tikz}
\usetikzlibrary{shapes.geometric, arrows.meta, positioning, calc, fit, backgrounds}

%==================== TABLES ====================
\usepackage{booktabs, longtable, array, multirow}

%==================== LAYOUT ====================
\usepackage{geometry}
\geometry{a4paper, margin=1in}
\usepackage{hyperref}
\usepackage{xcolor}
\usepackage{enumitem}
\usepackage{caption}
\usepackage{subcaption}
\usepackage{float}

%==================== STYLE ====================
\tikzstyle{startstop} = [ellipse, draw, thick, fill=blue!10, align=center, minimum width=3cm, minimum height=1cm]
\tikzstyle{process}   = [rectangle, draw, thick, fill=yellow!10, align=center, minimum width=3.5cm, minimum height=0.9cm]
\tikzstyle{decision}  = [diamond, draw, thick, fill=red!10, align=center, aspect=2.5, inner sep=1pt]
\tikzstyle{io}        = [trapezium, trapezium left angle=70, trapezium right angle=110, draw, thick, fill=green!10, align=center]
\tikzstyle{arrow}     = [->, thick, >=Stealth]

%==================== MATH COMMANDS ====================
\newcommand{\QOA}{Q^{OA}}
\newcommand{\QPLT}{Q^{PLT}}
\newcommand{\qOA}{q^{OA}}
\newcommand{\qPLT}{q^{PLT}}
\newcommand{\rOA}{r^{OA}}
\newcommand{\rPLT}{r^{PLT}}
\newcommand{\Inv}{I}
\newcommand{\pos}[1]{\left[#1\right]^+}
\newcommand{\ind}[1]{\mathbf{1}\!\left(#1\right)}

\renewcommand{\arraystretch}{1.4}

%==================== DOCUMENT ====================
\begin{document}

\begin{titlepage}
  \centering
  \vspace*{2cm}
  {\LARGE\bfseries Phương pháp Hybrid GA--ALNS\\[0.5em]
  cho bài toán Phân bổ và Chuyển ngang Hàng hóa\\[0.5em]
  Đa Nhà máy với Ràng buộc Tồn kho}\\[2cm]
  {\large Kết hợp Genetic Algorithm và\\
  Adaptive Large Neighborhood Search}\\[3cm]
  {\normalsize Tài liệu kỹ thuật -- Khóa luận tốt nghiệp}
  \vfill
  {\large \today}
\end{titlepage}

\tableofcontents
\newpage

%==================== SECTION 1 ====================

\section{Tổng quan phương pháp}

\subsection{Bài toán và động lực nghiên cứu}

Bài toán phân bổ và chuyển ngang hàng hóa đa nhà máy (Multi-Facility Allocation and Proactive Lateral Transshipment -- MFA-PLT) được mô tả bằng mô hình tối ưu hóa có cấu trúc hỗn hợp nguyên--thực (MILP). Mô hình bao gồm:

\begin{itemize}
  \item \textbf{Biến quyết định liên tục và nguyên:} lượng phân bổ OA ($\QOA_{p,i,t}$), số case đầy đủ ($\qOA_{p,i,t}$), phần dư ($\rOA_{p,i,t}$); tương tự cho PLT.
  \item \textbf{Ràng buộc phi tuyến tính tiềm ẩn:} hàm $\pos{\cdot}$ (phần dương) xuất hiện trong hàm mục tiêu.
  \item \textbf{Không gian nghiệm lớn:} số biến tăng theo $|P| \times |L|^2 \times |T|$.
  \item \textbf{Tính phụ thuộc thời gian:} tồn kho $I_{p,i,t}$ phụ thuộc vào quyết định trong quá khứ thông qua lead time $LT_i^{OA}$ và $LT_{i,j}^{PLT}$.
\end{itemize}

Những đặc điểm này khiến cho việc giải chính xác bằng MILP solver trở nên kém khả thi khi quy mô bài toán tăng lên. Do đó, một phương pháp meta-heuristic kết hợp được đề xuất.

\subsection{Kiến trúc tổng thể}

Phương pháp đề xuất kết hợp hai meta-heuristic:

\begin{enumerate}
  \item \textbf{Genetic Algorithm (GA):} đóng vai trò khám phá không gian nghiệm rộng (global exploration), duy trì sự đa dạng quần thể và tránh rơi vào cực tiểu địa phương.
  \item \textbf{Adaptive Large Neighborhood Search (ALNS):} đóng vai trò khai thác cục bộ (local exploitation), tinh chỉnh từng nghiệm cá thể thông qua các toán tử phá hủy--tái tạo có trọng số thích nghi.
\end{enumerate}

\medskip

\noindent\textbf{Cơ chế tích hợp:} ALNS được nhúng vào GA như một \emph{bước local search} ngay sau bước giải mã (decode). Cụ thể, sau khi mỗi nhiễm sắc thể được giải mã thành nghiệm thực, ALNS được kích hoạt để cải thiện nghiệm đó trước khi đánh giá fitness. Cơ chế này còn được gọi là \emph{Lamarckian evolution} vì kinh nghiệm học được trong local search được phản ánh ngược lại vào nhiễm sắc thể.

\subsection{Phân rã bài toán theo sản phẩm}

Vì ràng buộc năng lực $\sum_{i \in L} \QOA_{p,i,t} = CAP_{p,t}$ tách rời theo từng sản phẩm $p$, bài toán có thể được phân rã thành $|P|$ bài toán con độc lập, mỗi bài toán con tối ưu hóa phân bổ cho một sản phẩm. GA-ALNS được áp dụng song song hoặc tuần tự cho từng sản phẩm.

\begin{figure}[H]
\centering
\begin{tikzpicture}[
  node distance = 1.6cm,
  every node/.style = {font=\small}
]
  \node (p1) [process, minimum width=5cm] {Sản phẩm $p=1$: GA-ALNS};
  \node (p2) [process, below of=p1, minimum width=5cm] {Sản phẩm $p=2$: GA-ALNS};
  \node (dots) [below of=p2, yshift=0.3cm] {$\vdots$};
  \node (pp) [process, below of=dots, yshift=0.3cm, minimum width=5cm] {Sản phẩm $p=|P|$: GA-ALNS};
  \node (agg) [io, below of=pp, minimum width=5cm] {Tổng hợp $Z = \sum_p Z_p$};

  \draw[arrow] (p1) -- (p2);
  \draw[arrow] (p2) -- (dots);
  \draw[arrow] (dots) -- (pp);
  \draw[arrow] (pp) -- (agg);
\end{tikzpicture}
\caption{Phân rã bài toán theo từng sản phẩm}
\end{figure}

\subsection{Biểu diễn nhiễm sắc thể}

\subsubsection{Phiên bản biểu diễn}

Đối với mỗi sản phẩm $p$ đã được cố định, một nhiễm sắc thể $\mathcal{C}$ được định nghĩa như một ma trận:

\begin{equation}
  \mathcal{C} = \left(\QOA_{i,t}\right)_{i \in L,\, t \in T}
  \in \mathbb{Z}_+^{|L| \times |T|}
\end{equation}

trong đó $\QOA_{i,t}$ là \emph{tổng lượng} hàng được phân bổ từ nguồn trung tâm đến nhà máy $i$ tại giai đoạn $t$ (bỏ chỉ số $p$ vì đã cố định). Mỗi ô trong ma trận là một số nguyên không âm.

\subsubsection{Ràng buộc trên nhiễm sắc thể}

\begin{equation}
  \sum_{i \in L} \QOA_{i,t} = CAP_t, \quad \forall t \in T
  \tag{Ràng buộc năng lực -- cứng}
\end{equation}

\begin{equation}
  \QOA_{i,t} \geq 0, \quad \forall i \in L, t \in T
\end{equation}

\subsubsection{Biến suy ra (không encode)}

Các biến sau được suy ra trong quá trình giải mã, không phải encode trực tiếp:

\begin{align}
  \qOA_{i,t} &= \left\lfloor \frac{\QOA_{i,t}}{CP_{i,t}} \right\rfloor
  \tag{số case đầy đủ}\\[4pt]
  \rOA_{i,t} &= \QOA_{i,t} - CP_{i,t} \cdot \qOA_{i,t}
  \tag{phần dư OA}\\[4pt]
  \QPLT_{i,j,t} &= f(\QOA, I)
  \tag{suy ra từ trạng thái tồn kho}\\[4pt]
  I_{i,t} &= g(\QOA, \QPLT, BI_i, \Delta I)
  \tag{cân bằng tồn kho}
\end{align}

\subsubsection{Ví dụ minh họa -- nhiễm sắc thể với 2 nhà máy, 3 giai đoạn}

Giả sử $CAP_t = 100$ đơn vị cho mọi $t$. Một nhiễm sắc thể hợp lệ:

\[
\mathcal{C} =
\begin{pmatrix}
  \QOA_{1,1} & \QOA_{1,2} & \QOA_{1,3} \\
  \QOA_{2,1} & \QOA_{2,2} & \QOA_{2,3}
\end{pmatrix}
=
\begin{pmatrix}
  40 & 55 & 30 \\
  60 & 45 & 70
\end{pmatrix}
\]

Kiểm tra: $40+60=100$, $55+45=100$, $30+70=100$ $\checkmark$

\subsection{Hàm giải mã (Decode)}

\label{subsec:decode}

Hàm giải mã chuyển đổi nhiễm sắc thể $\mathcal{C}$ thành một nghiệm hoàn chỉnh gồm $(\QOA, \QPLT, I)$:

\begin{algorithm}[H]
\caption{Giải mã nhiễm sắc thể (Decode)}
\begin{algorithmic}[1]
\Require Nhiễm sắc thể $\mathcal{C} = (\QOA_{i,t})$, tham số mô hình
\Ensure Nghiệm $(\QOA, \QPLT, I)$

\State \textbf{Bước 1: Tính phần dư OA}
\For{$i \in L$, $t \in T$}
  \State $\qOA_{i,t} \gets \lfloor \QOA_{i,t} / CP_{i,t} \rfloor$
  \State $\rOA_{i,t} \gets \QOA_{i,t} - CP_{i,t} \cdot \qOA_{i,t}$
\EndFor

\State \textbf{Bước 2: Mô phỏng tồn kho và xác định PLT}
\State $I_{i,0} \gets BI_i$ \quad cho mọi $i \in L$

\For{$t = 1$ đến $|T|$}

  \State \textit{-- Tính tồn kho tạm thời trước PLT --}
  \State $I^{temp}_{i,t} \gets I_{i,t-1} 
  + Q^{OA}_{i,t - LT_i^{OA}} 
  + \sum_{j \neq i} Q^{PLT}_{j,i,t - LT_{j,i}^{PLT}} 
  + \Delta I_{i,t}$
  
  \State (nếu chỉ số thời gian < 1 thì các giá trị tương ứng bằng 0)

  \State \textit{-- Xác định PLT theo heuristic --}
  \If{$t \le \min_{j \in L} LT_j^{OA} - 1$}
    \State $\text{need}_i = [L_{i,t} - I^{temp}_{i,t}]^+$
    \State $\text{supply}_i = [I^{temp}_{i,t} - L_{i,t}]^+$
    \State Thực hiện heuristic PLT
  \EndIf

  \State \textit{-- Cập nhật tồn kho sau PLT --}
  \State $I_{i,t} =
  I^{temp}_{i,t}
  - \sum_{j \neq i} Q^{PLT}_{i,j,t}$

\EndFor

\State \Return $(\QOA, \QPLT, I)$
\end{algorithmic}
\end{algorithm}

\begin{algorithm}[H]
\caption{Heuristic phân bổ PLT (trong bước giải mã)}
\label{alg:plt_heuristic}
\begin{algorithmic}[1]
\Require $I^{temp}_{i,t}$, $L_{i,t}$, $CP_{j,t}$ (case-pack của nhà máy nhận)
\Ensure $\QPLT_{i,j,t}$ cho mọi cặp $(i,j)$

\State Xác định tập nhà máy thiếu: $S^- = \{i : I^{temp}_{i,t} < L_{i,t}\}$
\State Xác định tập nhà máy dư: $S^+ = \{i : I^{temp}_{i,t} > L_{i,t}\}$
\State Xếp hạng $S^-$ theo mức thiếu hụt giảm dần: $\delta_i^- = L_{i,t} - I^{temp}_{i,t}$
\State Xếp hạng $S^+$ theo mức dư thừa giảm dần: $\delta_j^+ = I^{temp}_{j,t} - L_{j,t}$

\For{$j \in S^+$ (theo thứ tự)}
  \State $avail_j \gets \delta_j^+$
  \For{$i \in S^-$ (theo thứ tự)}
    \State $\text{need}_i \gets \delta_i^-$
    \State $\text{transfer} \gets \min(avail_j, \text{need}_i)$
    \State $k \gets \lfloor \text{transfer} / CP_{i,t} \rfloor$ \quad \textit{-- làm tròn theo case-pack nhà máy nhận --}
    \State $\QPLT_{j,i,t} \gets k \cdot CP_{i,t}$
    \State $\qPLT_{j,i,t} \gets k$; \quad $\rPLT_{j,i,t} \gets 0$ \quad \textit{-- nếu chỉ cho phép case đầy đủ --}
    \State $avail_j \gets avail_j - \QPLT_{j,i,t}$
    \State $\delta_i^- \gets \delta_i^- - \QPLT_{j,i,t}$
  \EndFor
\EndFor

\State \Return $\QPLT$
\end{algorithmic}
\end{algorithm}

\subsection{Hàm tính Fitness}

Fitness của nhiễm sắc thể được định nghĩa là giá trị hàm mục tiêu $Z$ sau khi giải mã:

\begin{align}
  f(\mathcal{C}) &= Z(\QOA, \QPLT, I) \notag\\
  &= \underbrace{\sum_{i,t} \left( C^o_{i,t} \pos{I_{i,t}-U_{i,t}} + C^s_{i,t} \pos{L_{i,t}-I_{i,t}} + C^b_{i,t} \pos{-I_{i,t}} \right)}_{\text{chi phí tồn kho}} \notag\\
  &\quad + \underbrace{\sum_{i,t} C^p_{i,t} \cdot \ind{\rOA_{i,t}>0}}_{\text{phạt OA không đủ case}} \notag\\
  &\quad + \underbrace{\sum_{i,j \neq i,\, t \in T_j^{PLT}} \left( C^p_{i,j,t} \cdot \ind{\rPLT_{i,j,t}>0} + d_{i,j} \cdot TC \cdot \ind{\QPLT_{i,j,t}>0} \right)}_{\text{chi phí PLT}}
  \label{eq:fitness}
\end{align}

Vì mục tiêu là \textbf{tối thiểu hóa}, fitness tốt hơn tương ứng với giá trị $f$ nhỏ hơn.


%==================== SECTION 2 ====================

\section{Genetic Algorithm (GA)}

\subsection{Tổng quan GA trong bài toán MFA-PLT}

Genetic Algorithm là phương pháp tìm kiếm dựa trên quần thể, mô phỏng tiến hóa tự nhiên. Trong bài toán này, mỗi cá thể trong quần thể biểu diễn một phương án phân bổ $\QOA$ cho một sản phẩm. GA thực hiện các toán tử di truyền để khám phá không gian nghiệm rộng và dần hội tụ về vùng nghiệm tốt.

\subsection{Sơ đồ luồng GA}

\begin{figure}[H]
\centering
\begin{tikzpicture}[node distance=1.5cm, every node/.style={font=\small}]

  \node (start)     [startstop]                          {Bắt đầu};
  \node (init)      [process, below of=start]            {Khởi tạo quần thể $\mathcal{P}$\\(Algorithm 1)};
  \node (eval)      [process, below of=init]             {Đánh giá fitness\\$f(\mathcal{C})$ cho mọi $\mathcal{C} \in \mathcal{P}$};
  \node (term)      [decision, below of=eval, yshift=-0.4cm] {Đạt điều kiện\\dừng?};
  \node (output)    [startstop, below of=term, yshift=-0.4cm] {Xuất nghiệm tốt nhất};
  \node (select)    [process, right=3.5cm of eval]       {Selection\\(Tournament)};
  \node (cross)     [process, below of=select]           {Crossover\\(BLX-$\alpha$, Algorithm 3)};
  \node (mut)       [process, below of=cross]            {Mutation\\(Swap Alloc., Algorithm 4)};
  \node (decode)    [process, below of=mut]              {Decode + ALNS\\(Algorithms 5--8)};
  \node (newpop)    [process, below of=decode]           {Cập nhật quần thể\\($\mu + \lambda$ hoặc $\mu, \lambda$)};

  \draw[arrow] (start) -- (init);
  \draw[arrow] (init) -- (eval);
  \draw[arrow] (eval) -- (term);
  \draw[arrow] (term) -- node[left]{Có} (output);
  \draw[arrow] (term.east) -- ++(0.6,0) |- (select.west) node[pos=0.1, above]{Không};
  \draw[arrow] (select) -- (cross);
  \draw[arrow] (cross) -- (mut);
  \draw[arrow] (mut) -- (decode);
  \draw[arrow] (decode) -- (newpop);
  \draw[arrow] (newpop.west) -- ++(-0.6,0) |- (eval.east);

\end{tikzpicture}
\caption{Sơ đồ luồng Genetic Algorithm (GA)}
\label{fig:ga_flowchart}
\end{figure}

\subsection{Khởi tạo quần thể}

\begin{algorithm}[H]
\caption{Khởi tạo quần thể đa dạng}
\label{alg:init}
\begin{algorithmic}[1]
\Require $n_{pop}$, $L$, $T$, $CAP_t$, $BI_i$, $L_{i,t}$
\Ensure $\mathcal{P} = \{\mathcal{C}^{(1)}, \ldots, \mathcal{C}^{(n_{pop})}\}$

\State $\mathcal{P} \gets \emptyset$

\State \textit{-- Cá thể 1: Giải MILP để lấy nghiệm ưu tiên --}
\State Giải MILP nhỏ để thu $\QOA^*_{i,t}$; \; $\mathcal{C}^{(1)} \gets \QOA^*$
\State $\mathcal{P} \gets \mathcal{P} \cup \{\mathcal{C}^{(1)}\}$

\State \textit{-- Cá thể 2 đến $\lfloor 0.5 n_{pop} \rfloor$: Nhiễu xung quanh nghiệm MILP --}
\For{$k = 2$ \textbf{đến} $\lfloor 0.5 n_{pop} \rfloor$}
  \For{$i \in L$, $t \in T$}
    \State $\epsilon_{i,t} \sim \mathcal{U}(-\alpha, \alpha)$, với $\alpha = 0.15$
    \State $\widetilde{\QOA}_{i,t} \gets \max\!\left(0,\; \QOA^*_{i,t} \cdot (1 + \epsilon_{i,t})\right)$
  \EndFor
  \State Chuẩn hóa: $\QOA^{(k)}_{i,t} \gets \left\lfloor \dfrac{\widetilde{\QOA}_{i,t}}{\sum_{i'}\widetilde{\QOA}_{i',t}} \cdot CAP_t \right\rceil$, sửa tổng về $CAP_t$
  \State $\mathcal{P} \gets \mathcal{P} \cup \{\mathcal{C}^{(k)}\}$
\EndFor

\State \textit{-- Cá thể $\lfloor 0.5 n_{pop} \rfloor + 1$ đến $\lfloor 0.8 n_{pop} \rfloor$: Heuristic nhu cầu --}
\For{$k$ trong khoảng này}
  \For{$t \in T$}
    \State $\delta_{i,t} \gets \max(0,\; L_{i,t} - BI_i)$ \quad \textit{-- nhu cầu bù thiếu hụt --}
    \State $\rho_{i,t} \gets \delta_{i,t} / \sum_{i'} \delta_{i',t}$ \quad (nếu $\sum = 0$ thì $\rho_{i,t} = 1/|L|$)
    \State $\QOA^{(k)}_{i,t} \gets \lfloor \rho_{i,t} \cdot CAP_t \rceil$, sửa tổng về $CAP_t$
  \EndFor
  \State $\mathcal{P} \gets \mathcal{P} \cup \{\mathcal{C}^{(k)}\}$
\EndFor

\State \textit{-- Cá thể còn lại: Ngẫu nhiên Dirichlet --}
\For{$k = \lfloor 0.8n_{pop} \rfloor + 1$ \textbf{đến} $n_{pop}$}
  \For{$t \in T$}
    \State $(u_1,\ldots,u_{|L|}) \sim \mathrm{Dirichlet}(\mathbf{1})$
    \State $\QOA^{(k)}_{i,t} \gets \lfloor u_i \cdot CAP_t \rceil$, sửa tổng về $CAP_t$
  \EndFor
  \State $\mathcal{P} \gets \mathcal{P} \cup \{\mathcal{C}^{(k)}\}$
\EndFor

\State \Return $\mathcal{P}$
\end{algorithmic}
\end{algorithm}

\subsubsection*{Giải thích chi tiết và ví dụ}

\paragraph{Bước chuẩn hóa:} Sau khi làm tròn, tổng $\sum_i \QOA_{i,t}$ có thể lệch khỏi $CAP_t$ do sai số làm tròn. Ta sửa bằng cách điều chỉnh thêm/bớt 1 đơn vị tại nhà máy có phần thập phân lớn nhất (hoặc nhỏ nhất):

\vspace{0.3em}
\noindent\textbf{Ví dụ:} $|L|=2$, $CAP_t = 100$, $\widetilde{\QOA}_{1,t}=43.7$, $\widetilde{\QOA}_{2,t}=58.1$. Sau làm tròn: $44+58=102 \neq 100$. Giảm 2 đơn vị tại nhà máy có phần dư nhỏ nhất: nhà máy 1 (phần dư $0.7 > 0.1$, giảm nhà máy 2 trước). Kết quả: $44 + 56 = 100$.

\subsection{Cơ chế chọn lọc (Selection)}

\begin{algorithm}[H]
\caption{Tournament Selection}
\label{alg:selection}
\begin{algorithmic}[1]
\Require $\mathcal{P}$, kích thước tournament $k_{tour}$, số cá thể cần chọn $n_{select}$
\Ensure $\mathcal{P}^{select}$

\State $\mathcal{P}^{select} \gets \emptyset$
\While{$|\mathcal{P}^{select}| < n_{select}$}
  \State Chọn ngẫu nhiên $k_{tour}$ cá thể từ $\mathcal{P}$
  \State $\mathcal{C}^* \gets \arg\min_{\mathcal{C}} f(\mathcal{C})$ trong tập đã chọn
  \State $\mathcal{P}^{select} \gets \mathcal{P}^{select} \cup \{\mathcal{C}^*\}$
\EndWhile
\State \Return $\mathcal{P}^{select}$
\end{algorithmic}
\end{algorithm}

\noindent Tournament selection ưu tiên cá thể tốt hơn nhưng vẫn cho phép cá thể yếu hơn được chọn (khi $k_{tour}$ nhỏ), đảm bảo áp lực chọn lọc cân bằng. Tham số khuyến nghị: $k_{tour} = 3$.

\subsection{Lai ghép (Crossover) -- BLX-$\alpha$ có chuẩn hóa}

\begin{algorithm}[H]
\caption{BLX-$\alpha$ Crossover cho phân bổ hàng hóa}
\label{alg:crossover}
\begin{algorithmic}[1]
\Require $\mathcal{C}^{(1)}, \mathcal{C}^{(2)}$, xác suất lai ghép $p_c$, $L$, $T$, $CAP_t$
\Ensure $\mathcal{C}^{(1')}, \mathcal{C}^{(2')}$

\State $\mathcal{C}^{(1')} \gets \mathcal{C}^{(1)}$, \; $\mathcal{C}^{(2')} \gets \mathcal{C}^{(2)}$

\If{$\mathrm{rand}() > p_c$} \Return $\mathcal{C}^{(1')}, \mathcal{C}^{(2')}$
\EndIf

\State Chọn $\alpha \sim \mathcal{U}(0.3, 0.7)$

\For{$t \in T$}
  \For{$i \in L$}
    \State $\widetilde{Q}^{(1')}_{i,t} \gets \alpha \cdot \QOA^{(1)}_{i,t} + (1-\alpha) \cdot \QOA^{(2)}_{i,t}$
    \State $\widetilde{Q}^{(2')}_{i,t} \gets (1-\alpha) \cdot \QOA^{(1)}_{i,t} + \alpha \cdot \QOA^{(2)}_{i,t}$
  \EndFor
  \State \textit{-- Chuẩn hóa để đảm bảo $\sum_i = CAP_t$ --}
  \For{$s \in \{1', 2'\}$}
    \State $\widehat{Q}^{(s)}_{i,t} \gets \widetilde{Q}^{(s)}_{i,t} / \sum_{i'} \widetilde{Q}^{(s)}_{i',t} \cdot CAP_t$
    \State $\QOA^{(s)}_{i,t} \gets \lfloor \widehat{Q}^{(s)}_{i,t} \rceil$
  \EndFor
  \State Sửa tổng về chính xác $CAP_t$ cho cả $s = 1', 2'$
\EndFor

\State \Return $\mathcal{C}^{(1')}, \mathcal{C}^{(2')}$
\end{algorithmic}
\end{algorithm}

\subsubsection*{Ví dụ lai ghép}

$|L|=2$, $CAP_t = 100$, $\alpha = 0.6$:
\[
  \mathcal{C}^{(1)}: (40, 60) \quad \mathcal{C}^{(2)}: (70, 30)
\]
\[
  \widetilde{Q}^{(1')} = (0.6\cdot40 + 0.4\cdot70,\; 0.6\cdot60 + 0.4\cdot30) = (52, 48)
\]
\[
  \widetilde{Q}^{(2')} = (0.4\cdot40 + 0.6\cdot70,\; 0.4\cdot60 + 0.6\cdot30) = (58, 42)
\]
Cả hai đều có tổng $= 100$ $\checkmark$

\subsection{Đột biến (Mutation) -- Swap Allocation}

\begin{algorithm}[H]
\caption{Đột biến hoán đổi phân bổ (Swap Allocation Mutation)}
\label{alg:mutation}
\begin{algorithmic}[1]
\Require $\mathcal{C} = (\QOA_{i,t})$, xác suất đột biến $p_m$, $L$, $T$, $CAP_t$
\Ensure $\mathcal{C}'$

\State $\mathcal{C}' \gets \mathcal{C}$

\For{$t \in T$}
  \If{$\mathrm{rand}() \leq p_m$}
    \State Chọn ngẫu nhiên $i_1, i_2 \in L$, $i_1 \neq i_2$
    \State $\delta \sim \mathcal{U}_{\mathbb{Z}}\!\left(1,\; \min\!\left(\QOA_{i_1,t},\; \left\lfloor 0.2 \cdot CAP_t \right\rfloor\right)\right)$
    \State $\QOA'_{i_1,t} \gets \QOA_{i_1,t} - \delta$
    \State $\QOA'_{i_2,t} \gets \QOA_{i_2,t} + \delta$
    \State \textit{-- Ràng buộc $\sum_i \QOA'_{i,t} = CAP_t$ tự động thỏa mãn --}
  \EndIf
\EndFor

\State \Return $\mathcal{C}'$
\end{algorithmic}
\end{algorithm}

\subsubsection*{Ví dụ đột biến}

$|L|=3$, $CAP_t=100$, $(\QOA_{1,t}, \QOA_{2,t}, \QOA_{3,t}) = (30, 50, 20)$. Chọn $i_1=2$, $i_2=3$, $\delta=10$:
\[
  (30,\; 50-10,\; 20+10) = (30, 40, 30), \quad \text{tổng} = 100 \; \checkmark
\]

\subsection{Điều kiện dừng và cập nhật quần thể}

\begin{itemize}
  \item \textbf{Điều kiện dừng:} (1) đạt số thế hệ tối đa $G_{max}$; hoặc (2) không cải thiện fitness tốt nhất trong $G_{stag}$ thế hệ liên tiếp.
  \item \textbf{Cập nhật quần thể ($\mu + \lambda$):} Hợp nhất quần thể cha ($\mu$) và con ($\lambda$), giữ lại $n_{pop}$ cá thể tốt nhất. Đảm bảo elitism: cá thể tốt nhất luôn được giữ lại.
  \item \textbf{Loại bỏ trùng lặp:} Nếu hai cá thể giống hệt nhau (theo ma trận $\QOA$), thay thế bằng cá thể ngẫu nhiên mới để duy trì đa dạng quần thể.
\end{itemize}

\subsection{Tham số GA khuyến nghị}

\begin{center}
\begin{tabular}{llll}
\toprule
\textbf{Tham số} & \textbf{Ký hiệu} & \textbf{Giá trị} & \textbf{Ghi chú} \\
\midrule
Kích thước quần thể & $n_{pop}$ & $50$--$100$ & Tùy quy mô bài toán \\
Số thế hệ tối đa & $G_{max}$ & $200$--$500$ & \\
Xác suất lai ghép & $p_c$ & $0.8$ & \\
Xác suất đột biến & $p_m$ & $0.1$--$0.2$ & Per-period \\
Kích thước tournament & $k_{tour}$ & $3$ & \\
Ngưỡng stagnation & $G_{stag}$ & $50$ & \\
\bottomrule
\end{tabular}
\end{center}


%==================== SECTION 3 ====================

\section{Adaptive Large Neighborhood Search (ALNS)}

\subsection{Tổng quan ALNS}

ALNS là một dạng mở rộng của Large Neighborhood Search (LNS), trong đó nhiều toán tử phá hủy (destroy) và tái tạo (repair) được duy trì song song, và trọng số của từng toán tử được cập nhật thích nghi dựa trên hiệu quả lịch sử của chúng. Trong bài toán MFA-PLT, ALNS đóng vai trò tìm kiếm cục bộ (local search) sau mỗi lần giải mã nhiễm sắc thể.

\subsection{Sơ đồ luồng ALNS}

\begin{figure}[H]
\centering
\begin{tikzpicture}[node distance=1.6cm, every node/.style={font=\small}]

  \node (init)    [startstop]                            {Nhận nghiệm $\mathcal{C}_{in}$\\ từ GA Decode};
  \node (setup)   [process, below of=init]               {Khởi tạo\\$x = x_{in}$, $x^* = x_{in}$\\trọng số $w_d$, $w_r$};
  \node (iter)    [decision, below of=setup, yshift=-0.4cm] {Hết số\\vòng ALNS?};
  \node (done)    [startstop, left=3cm of iter]          {Trả về $x^*$\\cho GA};
  \node (sel)     [process, below of=iter, yshift=-0.4cm] {Chọn destroy op $d$\\và repair op $r$\\theo xác suất $\propto w_d$, $w_r$};
  \node (destroy) [process, below of=sel]                {Destroy: phá hủy\\$q$ giai đoạn/nhà máy\\từ nghiệm $x$};
  \node (repair)  [process, below of=destroy]            {Repair: tái phân bổ\\các phần bị phá hủy\\theo $r$};
  \node (decode2) [process, below of=repair]             {Decode + tính\\fitness $f(x')$};
  \node (accept)  [decision, below of=decode2, yshift=-0.3cm] {Accept\\$x'$?};
  \node (update)  [process, right=3.5cm of accept]       {Cập nhật\\trọng số $w_d$, $w_r$};
  \node (useX)    [process, above of=update]             {Cập nhật trạng thái:\\$x \gets x'$ (nếu accept)\\$x^* \gets x'$ (nếu best)};

  \draw[arrow] (init) -- (setup);
  \draw[arrow] (setup) -- (iter);
  \draw[arrow] (iter) -- node[above]{Có} (done);
  \draw[arrow] (iter) -- node[right]{Không} (sel);
  \draw[arrow] (sel) -- (destroy);
  \draw[arrow] (destroy) -- (repair);
  \draw[arrow] (repair) -- (decode2);
  \draw[arrow] (decode2) -- (accept);
  \draw[arrow] (accept) |- node[above left]{Là best} (useX);
  \draw[arrow] (accept.east) -- ++(0.3,0) |- node[right,pos=0.25]{Accept/Reject} (update.south);
  \draw[arrow] (useX) -- (update);
  \draw[arrow] (update.north) -- ++(0,0.5) -| (iter.east);

\end{tikzpicture}
\caption{Sơ đồ luồng ALNS}
\label{fig:alns_flowchart}
\end{figure}

\subsection{Không gian nghiệm trong ALNS}

Trong ALNS (được nhúng vào GA), nghiệm làm việc $x$ là nhiễm sắc thể đã giải mã, biểu diễn dưới dạng:

\[
  x = (\QOA_{i,t})_{i \in L,\, t \in T}
\]

ALNS thao tác trực tiếp trên ma trận $\QOA$ và đảm bảo ràng buộc $\sum_i \QOA_{i,t} = CAP_t$ sau mỗi bước repair.

\subsection{Toán tử Phá hủy (Destroy Operators)}

\subsubsection*{D1 -- Random Period Destroy}

Chọn ngẫu nhiên $q$ giai đoạn và xóa toàn bộ phân bổ trong các giai đoạn đó:

\begin{algorithm}[H]
\caption{D1: Random Period Destroy}
\label{alg:d1}
\begin{algorithmic}[1]
\Require Nghiệm $x = (\QOA_{i,t})$, số giai đoạn bị phá hủy $q$
\Ensure Nghiệm bị phá $x^-$, tập giai đoạn bị phá hủy $\mathcal{T}^-$

\State $\mathcal{T}^- \gets$ chọn ngẫu nhiên $q$ phần tử từ $T$
\State $x^- \gets x$
\For{$t \in \mathcal{T}^-$}
  \For{$i \in L$}
    \State $\QOA_{i,t}^- \gets 0$ \quad \textit{-- phần bị phá hủy sẽ được repair --}
  \EndFor
\EndFor
\State \Return $x^-$, $\mathcal{T}^-$

\end{algorithmic}
\end{algorithm}

\subsubsection*{D2 -- Worst Cost Destroy}

Xác định $q$ giai đoạn có đóng góp chi phí cao nhất và phá hủy chúng:

\begin{algorithm}[H]
\caption{D2: Worst Cost Destroy}
\label{alg:d2}
\begin{algorithmic}[1]
\Require Nghiệm $x$, số giai đoạn bị phá hủy $q$, trạng thái tồn kho $I$
\Ensure $x^-$, $\mathcal{T}^-$

\For{$t \in T$}
  \State $cost_t \gets \sum_{i \in L}\left( C^o_{i,t}\pos{I_{i,t}-U_{i,t}} + C^s_{i,t}\pos{L_{i,t}-I_{i,t}} + C^b_{i,t}\pos{-I_{i,t}} \right)$
\EndFor
\State Sắp xếp $T$ theo $cost_t$ giảm dần
\State $\mathcal{T}^- \gets$ $q$ phần tử đầu tiên
\State Đặt $\QOA_{i,t}^- \gets 0$ với mọi $i, t \in \mathcal{T}^-$
\State \Return $x^-$, $\mathcal{T}^-$

\end{algorithmic}
\end{algorithm}

\subsubsection*{D3 -- Facility-focused Destroy}

Xác định nhà máy có tồn kho lệch lớn nhất so với mức sàn/trần và phá hủy các giai đoạn liên quan:

\begin{algorithm}[H]
\caption{D3: Facility-focused Destroy}
\label{alg:d3}
\begin{algorithmic}[1]
\Require Nghiệm $x$, trạng thái $I$, số giai đoạn $q$
\Ensure $x^-$, $\mathcal{T}^-$

\For{$i \in L$}
  \State $\text{dev}_i \gets \sum_t \left(\pos{I_{i,t} - U_{i,t}} + \pos{L_{i,t} - I_{i,t}}\right)$
\EndFor
\State $i^* \gets \arg\max_{i} \text{dev}_i$
\State $\mathcal{T}^- \gets$ $q$ giai đoạn có $\pos{L_{i^*,t}-I_{i^*,t}}$ lớn nhất
\State Đặt $\QOA_{i,t}^- \gets 0$ với mọi $i \in L$, $t \in \mathcal{T}^-$
\State \Return $x^-$, $\mathcal{T}^-$

\end{algorithmic}
\end{algorithm}

\noindent\textbf{Tham số phá hủy:} $q \sim \mathcal{U}_{\mathbb{Z}}(q_{\min}, q_{\max})$, khuyến nghị $q_{\min} = \lceil 0.1|T| \rceil$, $q_{\max} = \lceil 0.4|T| \rceil$.

\subsection{Toán tử Tái tạo (Repair Operators)}

\subsubsection*{R1 -- Greedy Repair theo nhu cầu thiếu hụt}

Phân bổ lại $CAP_t$ cho các giai đoạn bị phá hủy dựa trên mức thiếu hụt tồn kho dự kiến:

\begin{algorithm}[H]
\caption{R1: Greedy Demand-based Repair}
\label{alg:r1}
\begin{algorithmic}[1]
\Require $x^-$, $\mathcal{T}^-$, $I^{prev}_{i,t}$ (tồn kho từ nghiệm chưa phá hủy), $L_{i,t}$, $CAP_t$
\Ensure Nghiệm được sửa $x'$

\State $x' \gets x^-$
\For{$t \in \mathcal{T}^-$ (theo thứ tự tăng dần)}
  \For{$i \in L$}
    \State $\delta_{i,t} \gets \max(0,\; L_{i,t} - I^{prev}_{i,t})$ \quad \textit{-- nhu cầu bù --}
  \EndFor
  \If{$\sum_i \delta_{i,t} = 0$}
    \State Phân bổ đều: $\QOA'_{i,t} \gets CAP_t / |L|$
  \Else
    \State $\QOA'_{i,t} \gets \delta_{i,t} / \sum_{i'} \delta_{i',t} \cdot CAP_t$
  \EndIf
  \State Làm tròn nguyên, sửa tổng về $CAP_t$
\EndFor
\State \Return $x'$

\end{algorithmic}
\end{algorithm}

\subsubsection*{R2 -- Random Repair}

\begin{algorithm}[H]
\caption{R2: Random Repair (Dirichlet)}
\label{alg:r2}
\begin{algorithmic}[1]
\Require $x^-$, $\mathcal{T}^-$, $CAP_t$
\Ensure $x'$

\State $x' \gets x^-$
\For{$t \in \mathcal{T}^-$}
  \State $(u_1,\ldots,u_{|L|}) \sim \mathrm{Dirichlet}(\mathbf{1})$
  \State $\QOA'_{i,t} \gets \lfloor u_i \cdot CAP_t \rceil$, sửa tổng về $CAP_t$
\EndFor
\State \Return $x'$

\end{algorithmic}
\end{algorithm}

\subsubsection*{R3 -- Regret-based Repair}

Tính \emph{regret} của mỗi nhà máy (sự chênh lệch chi phí nếu không được cấp phát) và ưu tiên nhà máy có regret lớn nhất:

\begin{algorithm}[H]
\caption{R3: Regret-based Repair}
\label{alg:r3}
\begin{algorithmic}[1]
\Require $x^-$, $\mathcal{T}^-$, $I^{prev}$, chi phí tồn kho
\Ensure $x'$

\State $x' \gets x^-$
\For{$t \in \mathcal{T}^-$}
  \For{$i \in L$}
    \State $\text{regret}_{i,t} \gets C^s_{i,t} \cdot \pos{L_{i,t} - I^{prev}_{i,t}} + C^b_{i,t} \cdot \pos{-I^{prev}_{i,t}}$
  \EndFor
  \State $\rho_{i,t} \gets \text{regret}_{i,t} / \sum_{i'} \text{regret}_{i',t}$
  \State $\QOA'_{i,t} \gets \lfloor \rho_{i,t} \cdot CAP_t \rceil$, sửa tổng về $CAP_t$
\EndFor
\State \Return $x'$

\end{algorithmic}
\end{algorithm}

\subsection{Tiêu chí chấp nhận (Acceptance Criteria)}

ALNS sử dụng tiêu chí chấp nhận dựa trên Simulated Annealing (SA):

\begin{equation}
  P(\text{accept}) =
  \begin{cases}
    1, & \text{nếu } f(x') \leq f(x) \\
    \exp\!\left(-\dfrac{f(x') - f(x)}{\theta}\right), & \text{ngược lại}
  \end{cases}
  \label{eq:acceptance}
\end{equation}

trong đó $\theta$ là nhiệt độ, giảm dần theo lịch làm nguội: $\theta \gets \theta \cdot \eta$ sau mỗi vòng lặp ($\eta = 0.99$). Nhiệt độ ban đầu $\theta_0$ được hiệu chỉnh sao cho $P(\text{accept ngẫu nhiên}) \approx 0.05$ tại đầu mỗi ALNS call.

\subsection{Cập nhật trọng số thích nghi (Adaptive Weight Update)}

\subsubsection*{Cơ chế tính điểm}

Sau mỗi vòng lặp ALNS, toán tử được chọn nhận điểm thưởng $\sigma$ dựa trên kết quả:

\begin{center}
\begin{tabular}{lll}
\toprule
\textbf{Kết quả} & \textbf{Điều kiện} & \textbf{Điểm $\sigma$} \\
\midrule
Nghiệm tốt toàn cục mới & $f(x') < f(x^*)$ & $\sigma_1 = 3$ \\
Nghiệm chấp nhận tốt & $f(x') < f(x)$ & $\sigma_2 = 2$ \\
Nghiệm chấp nhận (SA) & $f(x') \geq f(x)$, accept & $\sigma_3 = 1$ \\
Nghiệm bị từ chối & Reject & $\sigma_4 = 0$ \\
\bottomrule
\end{tabular}
\end{center}

\subsubsection*{Cập nhật trọng số}

Trọng số của toán tử $d$ (hoặc $r$) được cập nhật theo công thức:

\begin{equation}
  w_d \gets (1 - \lambda_{\text{rho}}) \cdot w_d + \lambda_{\text{rho}} \cdot \sigma_d
  \label{eq:weight_update}
\end{equation}

trong đó $\lambda_{\text{rho}} \in [0,1]$ là tốc độ học (khuyến nghị $\lambda_{\text{rho}} = 0.15$). Số lần sử dụng toán tử được theo dõi để ổn định ước lượng: cập nhật chỉ xảy ra sau mỗi $\Delta_{seg}$ vòng lặp (ví dụ $\Delta_{seg} = 20$).

\begin{algorithm}[H]
\caption{Cập nhật trọng số thích nghi ALNS}
\label{alg:weight_update}
\begin{algorithmic}[1]
\Require Trọng số $\{w_d\}_{d=1}^{n_D}$, $\{w_r\}_{r=1}^{n_R}$; điểm tích lũy $\{\pi_d\}$, $\{\pi_r\}$; số lần dùng $\{n_d\}$, $\{n_r\}$; tốc độ $\lambda_{\text{rho}}$
\Ensure Trọng số mới

\For{$d = 1$ đến $n_D$}
  \If{$n_d > 0$}
    \State $w_d \gets (1-\lambda_{\text{rho}}) \cdot w_d + \lambda_{\text{rho}} \cdot (\pi_d / n_d)$
  \EndIf
  \State $\pi_d \gets 0$; \; $n_d \gets 0$ \quad \textit{-- reset segment --}
\EndFor

\For{$r = 1$ đến $n_R$}
  \If{$n_r > 0$}
    \State $w_r \gets (1-\lambda_{\text{rho}}) \cdot w_r + \lambda_{\text{rho}} \cdot (\pi_r / n_r)$
  \EndIf
  \State $\pi_r \gets 0$; \; $n_r \gets 0$
\EndFor

\State \Return $\{w_d\}$, $\{w_r\}$

\end{algorithmic}
\end{algorithm}

\noindent Xác suất chọn toán tử được tính bằng softmax trọng số:
\begin{equation}
  P(\text{chọn toán tử } d) = \frac{w_d}{\sum_{d'} w_{d'}}
\end{equation}

\subsection{Toán tử ALNS -- Bảng tóm tắt}

\begin{center}
\begin{tabular}{lll}
\toprule
\textbf{Loại} & \textbf{Ký hiệu} & \textbf{Mô tả} \\
\midrule
\multirow{3}{*}{Destroy} & D1 & Random Period Destroy \\
                          & D2 & Worst Cost Destroy \\
                          & D3 & Facility-focused Destroy \\
\midrule
\multirow{3}{*}{Repair}  & R1 & Greedy Demand-based Repair \\
                          & R2 & Random Repair (Dirichlet) \\
                          & R3 & Regret-based Repair \\
\bottomrule
\end{tabular}
\end{center}

\noindent Tất cả các toán tử repair đều đảm bảo ràng buộc cứng $\sum_i \QOA_{i,t} = CAP_t$ sau khi sửa chữa.


%==================== SECTION 4 ====================

\section{Giải thuật kết hợp Hybrid GA--ALNS}

\subsection{Triết lý kết hợp}

Phương pháp kết hợp GA--ALNS trong nghiên cứu này được thiết kế theo mô hình \textbf{Lamarckian Hybrid}: ALNS đóng vai trò tìm kiếm cục bộ (local search / intensification) được nhúng vào vòng lặp GA tại bước giải mã, trước khi đánh giá fitness. Điều này có nghĩa là:

\begin{enumerate}
  \item GA tạo ra các nhiễm sắc thể thông qua selection, crossover, mutation.
  \item Mỗi nhiễm sắc thể được giải mã thành nghiệm ($\QOA$, $\QPLT$, $I$).
  \item ALNS tinh chỉnh (refine) nghiệm đó bằng cách phá hủy--tái tạo.
  \item Nghiệm sau ALNS được \emph{cập nhật ngược lại} vào nhiễm sắc thể (Lamarckian learning).
  \item Fitness được tính dựa trên nghiệm đã qua ALNS.
\end{enumerate}

\subsection{Điều kiện kích hoạt ALNS}

Để tránh tốn kém tính toán không cần thiết, ALNS không được kích hoạt cho mọi cá thể ở mọi thế hệ. Các quy tắc kích hoạt:

\begin{itemize}
  \item \textbf{Theo thế hệ:} ALNS được kích hoạt mỗi $\Delta_G$ thế hệ (ví dụ $\Delta_G = 5$).
  \item \textbf{Theo cá thể:} Tại thế hệ ALNS, chỉ áp dụng cho top-$k$ cá thể tốt nhất hoặc toàn bộ nếu quần thể nhỏ.
  \item \textbf{Theo stagnation:} ALNS luôn được kích hoạt khi phát hiện quần thể không cải thiện trong $G_{stag}^{ALNS}$ thế hệ.
\end{itemize}

\subsection{Sơ đồ luồng Hybrid GA--ALNS tổng thể}

\begin{figure}[H]
\centering
\resizebox{\linewidth}{!}{%
\begin{tikzpicture}[node distance=1.4cm, every node/.style={font=\small}]

  %============ CỘT TRÁI: GA ============
  \node (start)    [startstop]                          {Bắt đầu};
  \node (init)     [process, below of=start]            {Khởi tạo quần thể\\$\mathcal{P}$ (Alg.~\ref{alg:init})};
  \node (evalf)    [process, below of=init]             {Đánh giá\\fitness toàn bộ $\mathcal{P}$};
  \node (term)     [decision, below of=evalf, yshift=-0.3cm] {Điều kiện\\dừng GA?};
  \node (out)      [startstop, below of=term, yshift=-0.3cm] {Xuất $\mathcal{C}^*$};
  \node (sel)      [process, below of=out]              {Selection (Alg.~\ref{alg:selection})};
  \node (xover)    [process, below of=sel]              {Crossover (Alg.~\ref{alg:crossover})};
  \node (mut)      [process, below of=xover]            {Mutation (Alg.~\ref{alg:mutation})};
  \node (newpop)   [process, below of=mut]              {Cập nhật $\mathcal{P}$\\($\mu+\lambda$ strategy)};

  %============ CỘT PHẢI: DECODE + ALNS ============
  \node (dec)      [process, right=4.5cm of evalf]      {Decode $\mathcal{C}$\\$\Rightarrow$ $(\QOA, \QPLT, I)$};
  \node (actv)     [decision, below of=dec, yshift=-0.3cm] {Kích hoạt\\ALNS?};
  \node (alns)     [process, below of=actv, yshift=-0.3cm, fill=orange!15] {ALNS Local Search\\(Algs.~\ref{alg:d1}--\ref{alg:r3}, Alg.~\ref{alg:weight_update})};
  \node (lama)     [process, below of=alns]             {Lamarck update:\\$\mathcal{C} \gets$ encode$(x^*)$};
  \node (fit)      [process, below of=lama]             {Tính fitness $f(\mathcal{C})$\\(Phương trình~\ref{eq:fitness})};

  %============ ĐƯỜNG NGANG ============
  \node (gaplt)    [process, right=4.5cm of dec, fill=green!10]  {Mô phỏng PLT\\(Alg.~\ref{alg:plt_heuristic})};
  \node (inv)      [process, below of=gaplt]            {Cập nhật tồn kho\\$I_{i,t}$ (cân bằng tồn kho)};

  %============ MŨI TÊN CHÍNH ============
  \draw[arrow] (start) -- (init);
  \draw[arrow] (init) -- (evalf);
  \draw[arrow] (evalf) -- (term);
  \draw[arrow] (term) -- node[left]{Có} (out);
  \draw[arrow] (term.west) -- ++(-1.2,0) |- (sel.west) node[pos=0.12, above]{Không};
  \draw[arrow] (sel) -- (xover);
  \draw[arrow] (xover) -- (mut);
  \draw[arrow] (mut.east) -- ++(0.5,0) |- (dec.west);
  \draw[arrow] (dec) -- (gaplt);
  \draw[arrow] (gaplt) -- (inv);
  \draw[arrow] (inv.south) -- ++(0,-0.4) -| (actv.east);
  \draw[arrow] (actv) -- node[right]{Có} (alns);
  \draw[arrow] (alns) -- (lama);
  \draw[arrow] (lama) -- (fit);
  \draw[arrow] (actv.west) -- ++(-0.5,0) |- node[left,pos=0.7]{Không} (fit.west);
  \draw[arrow] (fit.west) -- ++(-0.5,0) |- (newpop.east);
  \draw[arrow] (newpop.west) -- ++(-0.5,0) |- (evalf.west);

\end{tikzpicture}
}%
\caption{Sơ đồ luồng đầy đủ Hybrid GA--ALNS}
\label{fig:hybrid_flowchart}
\end{figure}

\subsection{Cơ chế Lamarckian Learning}

Sau khi ALNS tìm được nghiệm cải thiện $x^* = (\QOA^*_{i,t})$, nhiễm sắc thể được cập nhật ngược bằng cách thay thế trực tiếp:

\begin{equation}
  \mathcal{C}' \gets x^* = (\QOA^*_{i,t})_{i \in L,\, t \in T}
\end{equation}

Điều này đảm bảo kiến thức học được từ local search được kế thừa trong các thế hệ tiếp theo.

\subsection{Giải thuật Hybrid GA--ALNS hoàn chỉnh}

\begin{algorithm}[H]
\caption{Hybrid GA--ALNS}
\label{alg:hybrid}
\begin{algorithmic}[1]
\Require Tham số mô hình, $n_{pop}$, $G_{max}$, $G_{stag}$, $\Delta_G$, $n_{ALNS}$
\Ensure Nghiệm tốt nhất $\mathcal{C}^*$, chi phí $Z^*$

\State $\mathcal{P} \gets \text{InitPopulation}(n_{pop})$ \quad [Alg.~\ref{alg:init}]
\For{$\mathcal{C} \in \mathcal{P}$}
  \State $f(\mathcal{C}) \gets \text{Decode-and-Evaluate}(\mathcal{C})$ \quad [Algs.~\ref{alg:plt_heuristic}, Eq.~\eqref{eq:fitness}]
\EndFor
\State $\mathcal{C}^* \gets \arg\min_{\mathcal{C} \in \mathcal{P}} f(\mathcal{C})$; \; $Z^* \gets f(\mathcal{C}^*)$
\State $g_{no\_improve} \gets 0$; \; $g \gets 0$
\State Khởi tạo trọng số ALNS: $w_d = w_r = 1.0$

\While{$g < G_{max}$ \textbf{và} $g_{no\_improve} < G_{stag}$}
  \State $g \gets g + 1$
  \State $\mathcal{P}^{select} \gets \text{TournamentSelection}(\mathcal{P})$ \quad [Alg.~\ref{alg:selection}]
  \State $\mathcal{P}^{off} \gets $ Crossover$(\mathcal{P}^{select})$ \quad [Alg.~\ref{alg:crossover}]
  \State $\mathcal{P}^{off} \gets $ Mutation$(\mathcal{P}^{off})$ \quad [Alg.~\ref{alg:mutation}]

  \For{$\mathcal{C} \in \mathcal{P}^{off}$}
    \State $x \gets \text{Decode}(\mathcal{C})$ \quad \textit{-- bao gồm mô phỏng PLT --}
    
    \If{$g \bmod \Delta_G = 0$ \textbf{hoặc} $g_{no\_improve} \geq G_{stag}/2$}
      \State $x^*, w_d, w_r \gets \text{ALNS}(x, n_{ALNS}, w_d, w_r)$ \quad [Algs.~\ref{alg:d1}--\ref{alg:weight_update}]
      \State $\mathcal{C} \gets x^*$ \quad \textit{-- Lamarckian update --}
      \State $f(\mathcal{C}) \gets \text{EvaluateFitness}(x^*)$
    \Else
      \State $f(\mathcal{C}) \gets \text{EvaluateFitness}(x)$
    \EndIf
  \EndFor

  \State $\mathcal{P} \gets \text{UpdatePopulation}(\mathcal{P}, \mathcal{P}^{off})$ \quad \textit{-- $(\mu+\lambda)$ + elitism --}

  \If{$\min_{\mathcal{C} \in \mathcal{P}} f(\mathcal{C}) < Z^*$}
    \State $Z^* \gets \min_{\mathcal{C}} f(\mathcal{C})$; \; $\mathcal{C}^* \gets \arg\min_{\mathcal{C}} f(\mathcal{C})$
    \State $g_{no\_improve} \gets 0$
  \Else
    \State $g_{no\_improve} \gets g_{no\_improve} + 1$
  \EndIf
\EndWhile

\State \Return $\mathcal{C}^*$, $Z^*$

\end{algorithmic}
\end{algorithm}

\subsection{Xử lý ràng buộc -- Hard vs. Soft Constraints}

\begin{center}
\begin{tabular}{p{0.45\textwidth}p{0.15\textwidth}p{0.3\textwidth}}
\toprule
\textbf{Ràng buộc} & \textbf{Loại} & \textbf{Cơ chế xử lý} \\
\midrule
$\sum_i \QOA_{i,t} = CAP_t$ & Cứng & Encode trực tiếp + repair sau crossover/mutation \\
$\QOA_{i,t} \geq 0$ & Cứng & Clip khi tạo/đột biến \\
$\sum_{j \neq i} \QPLT_{i,j,t} \leq I_{i,t} - L_{i,t}$ & Cứng & Đảm bảo trong Alg.~\ref{alg:plt_heuristic} \\
$I_{i,t} \geq 0$ & Mềm & Đưa vào hàm fitness (chi phí backorder $C^b$) \\
$I_{i,t} \leq U_{i,t}$ & Mềm & Đưa vào hàm fitness (chi phí overstock $C^o$) \\
$I_{i,t} \geq L_{i,t}$ & Mềm & Đưa vào hàm fitness (chi phí shortage $C^s$) \\
Case-pack $\QOA \equiv 0 \pmod{CP}$ & Mềm & Phần dư $r^{OA}$ gây phạt $C^p$ \\
\bottomrule
\end{tabular}
\end{center}

\noindent Ràng buộc cứng phải thỏa mãn trong mọi cá thể hợp lệ; ràng buộc mềm được "tha thứ" nhưng gây chi phí trong hàm mục tiêu.


%==================== SECTION 5 ====================

\section{Ví dụ minh họa}

\subsection{Thiết lập bài toán}

Xét bài toán nhỏ với:
\begin{itemize}
  \item 1 sản phẩm ($|P| = 1$), bỏ chỉ số $p$
  \item 2 nhà máy: $L = \{1, 2\}$
  \item 3 giai đoạn: $T = \{1, 2, 3\}$
  \item Lead time OA: $LT_1^{OA} = LT_2^{OA} = 2$ (hàng OA đến sau 2 giai đoạn)
  \item Lead time PLT: $LT_{1,2}^{PLT} = LT_{2,1}^{PLT} = 1$ (PLT mất 1 giai đoạn)
  \item $T_j^{PLT} = \{t : t \leq LT_j^{OA} - 1\} = \{1\}$ (PLT chỉ ở giai đoạn 1)
\end{itemize}

\subsection{Dữ liệu tham số}

\begin{center}
\begin{tabular}{lccc}
\toprule
\textbf{Tham số} & \textbf{Giai đoạn 1} & \textbf{Giai đoạn 2} & \textbf{Giai đoạn 3} \\
\midrule
$CAP_t$ & 100 & 100 & 100 \\
$L_{1,t}$ (mức sàn NM1) & 40 & 40 & 40 \\
$L_{2,t}$ (mức sàn NM2) & 30 & 30 & 30 \\
$U_{1,t}$ (mức trần NM1) & 120 & 120 & 120 \\
$U_{2,t}$ (mức trần NM2) & 100 & 100 & 100 \\
$\Delta I_{1,t}$ (tiêu dùng NM1) & $-20$ & $-25$ & $-30$ \\
$\Delta I_{2,t}$ (tiêu dùng NM2) & $-15$ & $-20$ & $-25$ \\
\bottomrule
\end{tabular}
\end{center}

\begin{center}
\begin{tabular}{lcc}
\toprule
\textbf{Tham số} & \textbf{Nhà máy 1} & \textbf{Nhà máy 2} \\
\midrule
$BI_i$ (tồn kho ban đầu) & 50 & 35 \\
$CP_{i,t}$ (case-pack) & 10 & 5 \\
$C^s_{i,t}$ (shortage cost) & 5 & 5 \\
$C^b_{i,t}$ (backorder cost) & 10 & 10 \\
$C^o_{i,t}$ (overstock cost) & 1 & 1 \\
$C^p_{i,t}$ (phạt OA) & 2 & 2 \\
\end{tabular}
\end{center}

Dữ liệu PLT: $d_{1,2} = d_{2,1} = 50$km, $TC = 0.1$, $C^p_{i,j,t} = 3$.

\subsection{Ví dụ Khởi tạo và Giải mã}

\subsubsection*{Bước 1: Nhiễm sắc thể khởi tạo}

Giả sử nghiệm MILP cho: $\QOA_{1,t}^* = (60, 55, 50)$, $\QOA_{2,t}^* = (40, 45, 50)$.

Nhiễm sắc thể tương ứng:
\[
\mathcal{C} =
\begin{pmatrix}
60 & 55 & 50 \\
40 & 45 & 50
\end{pmatrix}
, \quad
\text{kiểm tra: } 60+40=100, \; 55+45=100, \; 50+50=100 \; \checkmark
\]

\subsubsection*{Bước 2: Tính case-pack và phần dư OA}

Với $CP_{1,t}=10$ và $CP_{2,t}=5$:

\begin{center}
\begin{tabular}{ccccccc}
\toprule
& \multicolumn{2}{c}{$t=1$} & \multicolumn{2}{c}{$t=2$} & \multicolumn{2}{c}{$t=3$} \\
\cmidrule(lr){2-3}\cmidrule(lr){4-5}\cmidrule(lr){6-7}
& NM1 & NM2 & NM1 & NM2 & NM1 & NM2 \\
\midrule
$\QOA_{i,t}$ & 60 & 40 & 55 & 45 & 50 & 50 \\
$\qOA_{i,t} = \lfloor Q/CP \rfloor$ & 6 & 8 & 5 & 9 & 5 & 10 \\
$\rOA_{i,t} = Q - CP \cdot q$ & 0 & 0 & 5 & 0 & 0 & 0 \\
\bottomrule
\end{tabular}
\end{center}

Tại giai đoạn 2, NM1 có phần dư $\rOA_{1,2} = 5 \neq 0$ $\Rightarrow$ phạt $C^p_{1,2} = 2$.

\subsubsection*{Bước 3: Mô phỏng tồn kho giai đoạn 1}

$LT_i^{OA} = 2$ nên tại $t=1$: $\QOA_{i, 1-2} = \QOA_{i,-1} = 0$ (không có hàng OA nào đến).

\[
I^{temp}_{1,1} = BI_1 + 0 + \Delta I_{1,1} = 50 + 0 + (-20) = 30
\]
\[
I^{temp}_{2,1} = BI_2 + 0 + \Delta I_{2,1} = 35 + 0 + (-15) = 20
\]

Xét PLT tại $t=1 \in T_j^{PLT}$:
\begin{align*}
  &S^- = \{1, 2\}: \quad \delta_1^- = L_{1,1} - 30 = 40 - 30 = 10, \quad \delta_2^- = 30 - 20 = 10\\
  &S^+ = \emptyset \quad \text{(không có nhà máy nào dư)}
\end{align*}

Vì không có nhà máy dư, $\QPLT_{i,j,1} = 0$ cho mọi $i,j$.

\[
I_{1,1} = 30 \quad I_{2,1} = 20 \quad \text{(sau PLT)}
\]

Chi phí giai đoạn 1:
\[
C^s_{1,1} \pos{L_{1,1} - I_{1,1}} + C^s_{2,1} \pos{L_{2,1} - I_{2,1}} = 5 \cdot 10 + 5 \cdot 10 = 100
\]

\subsubsection*{Bước 4: Mô phỏng tồn kho giai đoạn 2}

$\QOA_{i,t-LT_i^{OA}} = \QOA_{i,2-2} = \QOA_{i,0}$: hàng đặt tại giai đoạn 0 (trước horizon, giả sử = 0).

\[
I^{temp}_{1,2} = I_{1,1} + 0 + \Delta I_{1,2} = 30 - 25 = 5
\]
\[
I^{temp}_{2,2} = I_{2,1} + 0 + \Delta I_{2,2} = 20 - 20 = 0
\]

$t=2 \notin T_j^{PLT}$ nên không có PLT. $I_{1,2} = 5$, $I_{2,2} = 0$.

Chi phí giai đoạn 2:
\[
5\cdot\pos{40-5} + 5\cdot\pos{30-0} + 10\cdot\pos{-0} + 2 \cdot \ind{\rOA_{1,2} > 0}
= 5\cdot35 + 5\cdot30 + 0 + 2 = 327
\]

\subsubsection*{Bước 5: Mô phỏng tồn kho giai đoạn 3}

Hàng OA đặt tại $t=1$ đến tại $t=3$ ($t-LT = 3-2=1$):

\[
I^{temp}_{1,3} = I_{1,2} + \QOA_{1,1} + \Delta I_{1,3} = 5 + 60 - 30 = 35
\]
\[
I^{temp}_{2,3} = I_{2,2} + \QOA_{2,1} + \Delta I_{2,3} = 0 + 40 - 25 = 15
\]

Chi phí giai đoạn 3:
\[
5\cdot\pos{40-35} + 5\cdot\pos{30-15} = 5\cdot5 + 5\cdot15 = 25 + 75 = 100
\]

\subsubsection*{Tổng fitness ban đầu}

\[
f(\mathcal{C}) = 100 + 327 + 100 = 527
\]

\subsection{Ví dụ ALNS -- Worst Cost Destroy và Greedy Repair}

\subsubsection*{Destroy (D2)} Giai đoạn có chi phí cao nhất: $t=2$ (chi phí 327). Phá hủy $q=1$ giai đoạn.

\[
x^- : \quad \QOA_{1,2}^- = 0, \quad \QOA_{2,2}^- = 0
\]

\subsubsection*{Repair (R1)} Nhu cầu bù tại $t=2$:
\[
\delta_{1,2} = \max(0, L_{1,2} - I_{1,1}) = \max(0, 40 - 30) = 10
\]
\[
\delta_{2,2} = \max(0, L_{2,2} - I_{2,1}) = \max(0, 30 - 20) = 10
\]
\[
\rho_{1,2} = 10/(10+10) = 0.5, \quad \rho_{2,2} = 0.5
\]
\[
\QOA_{1,2}' = \lfloor 0.5 \times 100 \rceil = 50, \quad \QOA_{2,2}' = 50
\]

Nhiễm sắc thể sau repair:
\[
\mathcal{C}' =
\begin{pmatrix}
60 & 50 & 50 \\
40 & 50 & 50
\end{pmatrix}
\]

\subsubsection*{Đánh giá lại} Với $\QOA_{1,2}' = 50$: $\qOA_{1,2} = \lfloor 50/10 \rfloor = 5$, $\rOA_{1,2} = 0$ $\Rightarrow$ không còn phạt case-pack.

Tồn kho giai đoạn 2 (OA đặt tại $t=0$, không có): $I_{1,2} = 5$, $I_{2,2} = 0$ (không thay đổi).

Chi phí giai đoạn 2 sau repair:
\[
5\cdot35 + 5\cdot30 + 0 = 325 \quad \text{(giảm từ 327, tiết kiệm 2 đơn vị phạt case-pack)}
\]

$f(\mathcal{C}') = 100 + 325 + 100 = 525 < 527 = f(\mathcal{C})$ $\Rightarrow$ chấp nhận $\mathcal{C}'$.

\subsection{Ví dụ cập nhật trọng số ALNS}

Sau vòng lặp này, cặp toán tử (D2, R1) tìm được nghiệm cải thiện:
\[
\sigma_{D2} = 2 \quad (\text{nghiệm tốt hơn nghiệm hiện tại})
\]
Cập nhật:
\[
w_{D2} \gets (1 - 0.15) \cdot 1.0 + 0.15 \cdot 2 = 0.85 + 0.30 = 1.15
\]

Lần sau, D2 có xác suất được chọn cao hơn D1 và D3 (ban đầu đều $= 1.0$):
\[
P(\text{D2}) = \frac{1.15}{1.0 + 1.15 + 1.0} = \frac{1.15}{3.15} \approx 36.5\%
\]


%==================== SECTION 6 ====================

\section{Phân tích độ phức tạp}

\subsection{Độ phức tạp không gian}

\subsubsection*{Không gian nghiệm}

Không gian nghiệm cho một sản phẩm được định nghĩa bởi:
\[
\mathcal{S} = \left\{ (\QOA_{i,t}) \in \mathbb{Z}_+^{|L| \times |T|} \;\middle|\; \sum_{i \in L} \QOA_{i,t} = CAP_t, \; \forall t \in T \right\}
\]

Kích thước không gian nghiệm xấp xỉ (số lượng nghiệm nguyên dương thỏa mãn tổng):
\[
|\mathcal{S}| \approx \prod_{t \in T} \binom{CAP_t + |L| - 1}{|L| - 1} \sim \mathcal{O}\!\left(\frac{CAP^{|L|-1}}{(|L|-1)!}\right)^{|T|}
\]

Với $|L|=6$, $|T|=8$, $CAP=1000$: $|\mathcal{S}| \sim \mathcal{O}(10^{40})$, xác nhận sự cần thiết của meta-heuristic.

\subsubsection*{Bộ nhớ quần thể}

Bộ nhớ cần cho quần thể GA: $\mathcal{O}(n_{pop} \cdot |L| \cdot |T|)$.

\subsection{Độ phức tạp thời gian}

\subsubsection*{Mỗi lần giải mã (Decode)}

Giải mã một nhiễm sắc thể bao gồm:
\begin{itemize}
  \item Tính $\qOA, \rOA$: $\mathcal{O}(|L| \cdot |T|)$
  \item Mô phỏng tồn kho theo thời gian: $\mathcal{O}(|L|^2 \cdot |T|)$ (do PLT mỗi giai đoạn xét $|L|^2$ cặp)
  \item Tính fitness: $\mathcal{O}(|L| \cdot |T|)$
\end{itemize}

\textbf{Tổng mỗi decode:} $\mathcal{O}(|L|^2 \cdot |T|)$.

\subsubsection*{Mỗi vòng ALNS}

Mỗi vòng ALNS gồm Destroy + Repair + Decode:
\[
\mathcal{O}(|L| \cdot |T|) + \mathcal{O}(|L| \cdot |T|) + \mathcal{O}(|L|^2 \cdot |T|) = \mathcal{O}(|L|^2 \cdot |T|)
\]

Nếu ALNS chạy $n_{ALNS}$ vòng: $\mathcal{O}(n_{ALNS} \cdot |L|^2 \cdot |T|)$.

\subsubsection*{Mỗi thế hệ GA}

\begin{center}
\begin{tabular}{lll}
\toprule
\textbf{Bước} & \textbf{Chi phí} & \textbf{Ghi chú} \\
\midrule
Khởi tạo & $\mathcal{O}(n_{pop} \cdot |L|^2 \cdot |T|)$ & Chỉ lần đầu \\
Selection & $\mathcal{O}(n_{pop} \cdot k_{tour})$ & Tournament \\
Crossover & $\mathcal{O}(n_{pop} \cdot |L| \cdot |T|)$ & BLX-$\alpha$ \\
Mutation & $\mathcal{O}(n_{pop} \cdot |T|)$ & Swap per period \\
Decode & $\mathcal{O}(n_{pop} \cdot |L|^2 \cdot |T|)$ & Dominant term \\
ALNS & $\mathcal{O}(n_{pop} / \Delta_G \cdot n_{ALNS} \cdot |L|^2 \cdot |T|)$ & Mỗi $\Delta_G$ thế hệ \\
\bottomrule
\end{tabular}
\end{center}

\subsubsection*{Tổng độ phức tạp}

\[
\mathcal{O}\!\left(G_{max} \cdot n_{pop} \cdot \left(1 + \frac{n_{ALNS}}{\Delta_G}\right) \cdot |L|^2 \cdot |T|\right)
\]

Với tham số điển hình $G_{max}=200$, $n_{pop}=50$, $n_{ALNS}=50$, $\Delta_G=5$, $|L|=6$, $|T|=8$:
\[
200 \times 50 \times (1 + 10) \times 36 \times 8 \approx 3.2 \times 10^7 \text{ phép tính cơ bản}
\]

So với MILP (NP-hard), đây là mức có thể chấp nhận được cho thực thi trong thời gian thực.

\subsection{Phân tích hội tụ}

\subsubsection*{Đảm bảo tìm tối ưu toàn cục}

GA với mutation không bằng 0 và quần thể đủ lớn có thể tìm được tối ưu toàn cục với xác suất tiến tới 1 khi $G_{max} \to \infty$ (theo định lý hội tụ GA -- Holland, 1975). Tuy nhiên, trong thực tế, hội tụ chỉ được đảm bảo về tối ưu địa phương.

\subsubsection*{Vai trò của ALNS}

ALNS giúp mở rộng vùng lân cận được khám phá trong mỗi bước local search, tương đương với việc tăng "bước nhảy" trong không gian nghiệm mà không tăng kích thước quần thể. Kết hợp với GA giúp thoát khỏi cực tiểu địa phương thông qua cơ chế tiến hóa.

\subsection{So sánh các phương pháp}

\begin{center}
\begin{tabular}{p{0.22\textwidth}cccc}
\toprule
\textbf{Phương pháp} & \textbf{Optimal?} & \textbf{Scalable?} & \textbf{Diversity} & \textbf{Local expl.} \\
\midrule
MILP (Exact) & $\checkmark$ & $\times$ & N/A & $\checkmark$ \\
GA thuần & Không đảm bảo & $\checkmark$ & Cao & Thấp \\
ALNS thuần & Không đảm bảo & $\checkmark$ & Thấp & Cao \\
\textbf{Hybrid GA-ALNS} & Không đảm bảo & $\checkmark$ & \textbf{Trung bình} & \textbf{Cao} \\
\bottomrule
\end{tabular}
\end{center}


\end{document}
